\documentclass[10pt]{article}
\usepackage[T1]{fontenc}
\usepackage[utf8]{inputenc}
\usepackage[french]{babel}
\usepackage{amsmath,amssymb,mathrsfs,mathpazo}
\usepackage{geometry}
\geometry{a4paper,margin=2cm}
\usepackage{graphicx,array,tabularx,tabularray,colortbl,multirow,multicol,booktabs,enumitem,pifont,ulem,fancybox,diagbox,cancel}
\usepackage{pstricks,pstricks-add,pst-plot,pst-tree,pst-node,pst-text,pst-func,pst-math,pst-bspline,pst-3dplot,pst-eucl,pst-solides3d}
\let\pscalculate\undefined
\usepackage{pst-bezier}
\usepackage{tikz,pgfplots,tkz-tab}
\usetikzlibrary{arrows,arrows.meta,patterns,calc,positioning,trees,decorations.markings,shapes.geometric,shapes.misc,shapes.arrows,intersections,decorations.pathmorphing,decorations.pathreplacing,backgrounds,shadows,babel}
\pgfplotsset{compat=1.18}
\usepgfplotslibrary{fillbetween}
\usepackage[np]{numprint}
\usepackage{listings,makecell,pas-tableur,twoopt}
\let\barre\undefined
\usepackage[pstricks]{tablvar}
\definecolor{darkgreen}{RGB}{0,100,0}
\providecommand{\red}{\color{red}}
\providecommand{\blue}{\color{blue}}
\providecommand{\green}{\color{green}}
\providecommand{\black}{\color{black}}
\providecommand{\white}{\color{white}}
\providecommand{\gray}{\color{gray}}
\providecommand{\cyan}{\color{cyan}}
\providecommand{\magenta}{\color{magenta}}
\providecommand{\yellow}{\color{yellow}}
\providecommand{\violet}{\color{violet}}
\providecommand{\orange}{\color{orange}}
\providecommand{\euro}{\texteuro}
\providecommand{\EUR}{\texteuro}
\providecommand{\ieme}{\textsuperscript{e}}
\providecommand{\degres}{\ensuremath{^\circ}}
\providecommand{\degree}{\ensuremath{^\circ}}
\providecommand{\vv}[1]{\overrightarrow{#1}}
\providecommand{\Coord}[1]{#1}
\providecommand{\decosix}{}
\providecommand{\starredbullet}{\ensuremath{\bullet}}
\pagestyle{empty}
\setlength{\parindent}{0pt}
\usepackage{adjustbox}\DeclareUnicodeCharacter{2236}{:}\DeclareUnicodeCharacter{2190}{\ensuremath{\leftarrow}}
\providecommand{\R}{\mathbb{R}}
\providecommand{\N}{\mathbb{N}}
\providecommand{\Z}{\mathbb{Z}}
\providecommand{\Q}{\mathbb{Q}}
\providecommand{\D}{\mathbb{D}}
\providecommand{\C}{\mathbb{C}}
\providecommand{\e}{{\mathrm{e}}}
\providecommand{\vect}[1]{\overrightarrow{#1}}
\providecommand{\vectt}[1]{\overrightarrow{\mathrm{#1}}}
\providecommand{\vvt}[1]{\overrightarrow{\mathrm{#1}}}
\providecommand{\np}[2][]{{#2~#1}}
\providecommand{\vv}[1]{{\overrightarrow{#1}}}
\renewcommand{\d}{{\mathrm{d}}}
\providecommand{\degres}{{^\circ}}
\providecommand{\degre}{{^\circ}}
\providecommand{\blue}{\color{blue}}
\providecommand{\red}{\color{red}}
\providecommand{\enskip}{\quad}
\providecommand{\llbracket}{\lbrack\!\lbrack}
\providecommand{\rrbracket}{\rbrack\!\rbrack}
\providecommand{\euro}{{\text{€}}}
\providecommand{\no}{{\text{n°}}}
\providecommand{\interval}[2]{{[#1;#2]}}
\providecommand{\intervFO}[2]{{[#1;#2[}}
\providecommand{\dotfill}{\ldots}

\providecommand{\ds}{\displaystyle}
\providecommand{\textding}[1]{#1}
\providecommand{\Oij}{\(\mathrm{O};\vec{\imath},\vec{\jmath}\)}
\providecommand{\Oijk}{\(\mathrm{O};\vec{\imath},\vec{\jmath},\vec{k}\)}
\providecommand{\Ouv}{\(\mathrm{O};\vec{u},\vec{v}\)}
\providecommand{\intervFF}[2]{\ensuremath{[#1;#2]}}
\renewcommand{\euro}{\ifmmode\text{\texteuro}\else\texteuro\fi}
\DeclareUnicodeCharacter{00B0}{\ifmmode^{\circ}\else\textdegree\fi}
\DeclareUnicodeCharacter{20AC}{\ifmmode\text{\texteuro}\else\texteuro\fi}
% Énoncés : Éducation nationale. Transcription et corrigés : APMEP, Denis Vergès et contributeurs. Découpage et classement : CoopMaths, Sébastien Lozano, Éric Elter et contributeurs. Adaptation de la présentation : Maths Explorer.
% CC BY-SA : https://coopmaths.fr/www/about/
% https://coopmaths.fr/alea/static/e3c/2020/tex/e3c_2020_00_sujet26_1_cor.tex
\begin{document}
\noindent\textit{APMEP / CoopMaths. CC BY-SA. Présentation adaptée par Maths Explorer.}\par\medskip
\subsection*{Question 1}
	Pour tout réel \(x\),
	\begin{align*}
		\dfrac{e^{2x}}{e^x+1} &= e^{2x - x - 1} \\
		&= e^{x - 1}.
	\end{align*}
	
	\subsection*{Question 2}
	Les points \(M(x\,;\,y)\) communs aux deux courbes ont des abscisses qui vérifient :
	\begin{align*}
		&15x^2 + 10x - 1 = 19x^2 - 22x + 10 \\
		\iff & 4x^2 - 32x + 11=0.
	\end{align*}
	Or pour cette équation du second degré :
	\[
	\Delta = 32^2 - 4 \times 4 \times 11 = 1024 - 176 = 848 > 0.
	\]
	Le discriminant est supérieur à zéro, donc cette équation a deux solutions distinctes ; les deux courbes ont deux points communs.
	
	\subsection*{Question 3}
	\begin{align*}
		&M(x\,;\,y) \in \mathcal{C}(A, R = 5)\\
		\iff	&AM^2 = 5^2 \\
		\iff&   (x - 3)^2 + (y - (-1))^2 = 25 \\
		\iff&   (x - 3)^2 + (y + 1)^2 = 25.
	\end{align*}
	
	\subsection*{Question 4}
	Un vecteur directeur de \((d)\) est par exemple \(\vec{d}\begin{pmatrix}-2 \\ 3\end{pmatrix}\), donc un vecteur normal est par exemple \(\vec{n}\begin{pmatrix}3 \\ 2\end{pmatrix}.\)
	
	\subsection*{Question 5}
	On sait que \(1 + 2 + 3 + \cdots + n = \dfrac{n(n + 1)}{2}.\)
	
	Il faut donc trouver le plus petit naturel tel que :
	\begin{align*}
		&\dfrac{n(n + 1)}{2} > 5000 \\
		\iff \quad & n(n + 1) > 10000 \\
		\iff \quad & n^2 + n - 10000 > 0.
	\end{align*}
	
	Pour le trinôme \(n^2 + n - 10000\), on a 
	\[
	\Delta = 1 + 40000 = 40001 > 0.
	\]
	Le trinôme a donc deux racines :
	\[
	n_1 = \dfrac{-1 + \sqrt{40001}}{2} \approx 99,5 \quad \text{et} \quad n_2 = \dfrac{-1 - \sqrt{40001}}{2} \approx -100,5.
	\]
	On sait que le trinôme est négatif sur l'intervalle \([n_2\,;\,n_1]\), donc le plus petit entier pour lequel le trinôme est positif est \(100\).

\end{document}
